\documentclass[exercice]{thedocument}%
\startdatakeys
\source{}
\cycle{}
\competencesDuSocle{}
\connaissancesRequises{}
\competencesTravaillees{}
\enddatakeys
\begin{document}%
En utilisant l'équerre ou le compas, construire le symétrique des figures
suivantes par rapport $(d)$.\par\bigskip
\begin{tikzpicture}[rotate=-45]
    \coordinate(d1)at(0,0);
    \coordinate[label=below right:$(d)$](d2)at(0,5);
    \coordinate[label=above right:A](A)at(-2,3);
    \draw(d1)--(d2);
    \draw node at(A){+};
\end{tikzpicture}\hskip150pt
\begin{tikzpicture}[rotate=-15]
    \coordinate(d1)at(0,0);
    \coordinate[label=below right:$(d)$](d2)at(0,5);
    \draw(d1)--(d2);
    \draw (-1,1)--(-3,3)--(-2,4)--cycle;
\end{tikzpicture}\par
\begin{tikzpicture}[rotate=20]
    \coordinate(d1)at(0,0);
    \coordinate[label=below right:$(d)$](d2)at(0,5);
    \coordinate(A)at(-3,2);
    \draw(d1)--(d2);
    \draw node at(A){+};
    \draw (A)circle(2);
\end{tikzpicture}\hskip150pt
\begin{tikzpicture}[rotate=20]
    \coordinate(d1)at(0,0);
    \coordinate[label=below right:$(d)$](d2)at(0,5);
    \coordinate(e1)at(3,2);
    \coordinate(e2)at(-1,4);
    \draw(d1)--(d2);
    \draw(e1)--(e2);
\end{tikzpicture}
\par
\end{document}
